\documentclass[10pt]{article}
\usepackage[margin=0.95in]{geometry}
\usepackage{amsmath,amssymb,amsthm,microtype}
\usepackage[hidelinks]{hyperref}
\newtheorem{theorem}{Theorem}
\newtheorem{lemma}[theorem]{Lemma}
\newcommand{\TT}{\mathcal T}
\newcommand{\Ext}{\operatorname{Ext}}
\title{Original-ball paths in Teichm\"uller space:\\
an elementary corollary of Lenzhen--Rafi}
\author{Alec Kriebel\thanks{Independent researcher. ORCID:
\href{https://orcid.org/0009-0001-9320-500X}{0009-0001-9320-500X}.}}
\date{13 October 2026\quad Version 1.0\quad Unrefereed expository note}
\begin{document}
\maketitle
\begin{abstract}
We record an elementary path corollary of Lenzhen--Rafi's established
uniform quasiconvexity theorem. For each fixed closed oriented surface of
genus at least two, every pair of points in any closed Teichm\"uller ball
can be joined inside that same ball by a continuous rectifiable path whose
length is at most their ambient distance plus a topology-dependent
constant. Radially moving the endpoints inward before joining them gives
the deduction. This supplies an affirmative derivation of Farb's
Question~3.5 from prior Teichm\"uller geometry. We make no claim to
historical priority or a new substantive Teichm\"uller theorem.
\end{abstract}

\section{The original-ball conclusion}
Fix a closed oriented surface $S$ of genus $g\geq2$. Its marked
Teichm\"uller space $\TT(S)$ has the metric
$d_T(X,Y)=\tfrac12\log K(X,Y)$, where $K(X,Y)$ is the least
quasiconformal dilatation in the marking class. Balls below are closed
balls for this metric, and lengths are ordinary metric lengths of paths.

\begin{theorem}\label{thm:main}
There is a finite constant $c_g\geq0$ such that, for every center $o$,
every finite radius $R\geq0$, and every $x,y\in B_{d_T}(o,R)$, there is
a continuous rectifiable path in $B_{d_T}(o,R)$ joining $x$ to $y$ with
length at most $d_T(x,y)+4c_g$.
\end{theorem}

The same $c_g$ works for all centers, radii and endpoints at the fixed
genus. In particular, ambient distance at most $k$ gives a path of length
at most $N_g(k)=k+4c_g$. Taking $k=2$ satisfies Farb's definition of
almost convexity and Question~3.5, which asks for bounded-length paths
inside the original radius ball~\cite[pp.~24--25]{Farb}.
No numerical value, genus-uniform bound or optimal constant is asserted.

The substantive geometric input is Lenzhen--Rafi's published uniform ball
theorem~\cite[Theorems~C and~17]{LR}. It allows a minimizing segment to
leave the endpoint-radius ball by a bounded additive amount. The short
general deduction below converts that allowance into a path in the
original ball. It requires no thickness condition and no replacement of
the Teichm\"uller metric.

\section{The elementary metric deduction}
\begin{lemma}\label{lem:retreat}
Let $(M,d)$ be a geodesic metric space. Suppose one finite $c\geq0$
has the following property: for every $o,a,b\in M$, some minimizing
segment from $a$ to $b$ satisfies
\begin{equation}\label{eq:radial}
 d(o,z)\leq\max\{d(o,a),d(o,b)\}+c
 \quad\text{at every point }z\text{ on that segment}.
\end{equation}
Then every $x,y\in B(o,R)$ can be joined by a continuous rectifiable
path in $B(o,R)$ of length at most $d(x,y)+4c$.
\end{lemma}
\begin{proof}
If $R<c$, join $x$ to $o$ and $o$ to $y$ by radial minimizing segments.
Their concatenation lies in $B(o,R)$ and has length at most $2R$, which
is at most $d(x,y)+4c$.

Suppose $R\geq c$. On radial minimizing segments from $x$ and $y$
toward $o$, move inward by
\[
 \alpha=\min\{c,d(o,x)\},\qquad
 \beta=\min\{c,d(o,y)\},
\]
obtaining $a$ and $b$. If an original endpoint had radial distance below
$c$, its new point is $o$; otherwise its radius has decreased by $c$.
Thus $d(o,a),d(o,b)\leq R-c$. Choose the $a$-to-$b$ segment in
\eqref{eq:radial}; all of it lies in $B(o,R)$. The two radial end pieces
also lie in that ball. The resulting concatenation has length
\[
 \alpha+d(a,b)+\beta
 \leq d(x,y)+2\alpha+2\beta
 \leq d(x,y)+4c.
\]
Arclength parametrizations of the three pieces paste to give an actual
continuous rectifiable path. Constant pieces cover $a=b$, coincident
endpoints and endpoints at the center. The case $R=c$ is included; so
are $c=0$ and $R=0$. No attainment of an intrinsic-distance infimum,
properness, or continuous choice of segments is needed.
\end{proof}

\section{Application of the established theorem}
Lenzhen--Rafi's extremal-length result~\cite[Theorem~15]{LR} gives a
finite $K_g\geq1$, depending only on the topology, such that
\begin{equation}\label{eq:ext}
 \Ext_z(\gamma)\leq K_g
 \max\{\Ext_a(\gamma),\Ext_b(\gamma)\}
\end{equation}
for every essential simple closed curve $\gamma$ and every point $z$
on the minimizing Teichm\"uller segment from $a$ to $b$. The constant
is uniform in the curve, endpoints and intermediate point. Their
Theorem~17 proves the corresponding radial ball statement. For clarity,
it can also be obtained directly from Kerckhoff's formula~\cite{Kerckhoff},
in the normalization recorded in~\cite[Section~2, equation~(2.2)]{FR}:
\[
 d_T(o,w)=\frac12\sup_\gamma
 \log\frac{\Ext_w(\gamma)}{\Ext_o(\gamma)}.
\]
Divide \eqref{eq:ext} by the positive $\Ext_o(\gamma)$ and take the
supremum of half logarithms. This gives
\[
 d_T(o,z)\leq
 \max\{d_T(o,a),d_T(o,b)\}+\frac12\log K_g.
\]
There is no need to select a curve or foliation attaining the supremum.
Minimizing Teichm\"uller segments exist between every pair of points;
see also~\cite[Section~2]{FR}. Lemma~\ref{lem:retreat}, with
$c_g=\tfrac12\log K_g$, proves Theorem~\ref{thm:main}.

The original published input and these geodesic and distance-formula
facts are established external results. This note does not reprove their
foundational Teichm\"uller theory. The journal's extremal-length and ball
theorem numbers are respectively 15 and 17; the earlier
arXiv:1002.3939v1 uses 13 and 15 for those body theorems.

\section{Attribution, scope and disclosure}
All substantive Teichm\"uller estimates here are due to Lenzhen--Rafi.
The contribution of this exposition is to write out the radial-retreat
implication with the exact original-ball and path-length quantifiers.
Historical priority for that elementary implication or this corollary
has not been established. The bounded source audit located no explicit
full-scope predecessor in the inspected bodies; it does not establish
worldwide absence or that Farb's question remained open in 2026.

Fortier Bourque--Rafi's nonconvexity result~\cite{FR} concerns the direct
minimizing segment and is compatible with the concatenated path above.
Petyt--Zalloum's Theorem~I~\cite{PZ}, present in their 2024 preprint and
2026 published version, gives almost convexity with $N(k)=k+6$ in a
dual wall space. Its map from the original space is quasiisometric; it
does not identify original Teichm\"uller balls and path lengths. We make
no comparison of numerical constants across those different metrics.
Neither the lemma nor this note asserts optimality of $4c_g$ in
Teichm\"uller space.

AI tools were used extensively in deriving the corollary, verifying the
proof, adversarial review, source auditing and preparation. This is an
unrefereed expository preprint. Its AI adversarial audits do not constitute
conventional human peer review, journal acceptance or formal proof-assistant
verification. Supporting authored audit records and source identities are
provided separately; third-party paper bodies are not redistributed.

\begin{thebibliography}{9}\small
\bibitem{Farb} B.~Farb, \emph{Some problems on mapping class groups and
moduli space}, in \emph{Problems on Mapping Class Groups and Related
Topics}, 2006; Question~3.5, pp.~24--25.
\url{https://www.math.uchicago.edu/~farb/papers/mcgbook.pdf}
\bibitem{LR} A.~Lenzhen and K.~Rafi, \emph{Length of a curve is
quasi-convex along a Teichm\"uller geodesic}, J. Differential Geom.
\textbf{88} (2011), 267--295; Theorems~C, 15 and 17.
\url{https://www.math.toronto.edu/rafi/Papers/Convexity.pdf}
\bibitem{Kerckhoff} S.~P.~Kerckhoff, \emph{The asymptotic geometry of
Teichm\"uller space}, Topology \textbf{19} (1980), 23--41.
\url{https://doi.org/10.1016/0040-9383(80)90029-4}
\bibitem{FR} M.~Fortier Bourque and K.~Rafi, \emph{Non-convex balls in
the Teichm\"uller metric}, J. Differential Geom. \textbf{110} (2018),
379--412; arXiv:1606.05170v3 (29 August 2017).
\url{https://arxiv.org/abs/1606.05170v3}
\bibitem{PZ} H.~Petyt and A.~Zalloum, \emph{Constructing metric spaces
from systems of walls}, Math. Ann. \textbf{396}, article 38 (2026);
Theorem~I. Earlier versions: arXiv:2404.12057v1 (2024) and v2 (2026).
\url{https://doi.org/10.1007/s00208-026-03534-1}
\end{thebibliography}
\end{document}
